\section{El cambio organizacional y la estrategia de DeepMind}

\subsection{Integración tres en uno: la construcción de un superequipo de IA}

La creación de Google DeepMind fue una integración organizacional a gran escala: la fusión de tres equipos —Google Brain, Google Research y el DeepMind original— en una sola entidad unificada. La lógica central de esta integración fue: \textbf{sustituir la dispersión propia de una gran empresa por el enfoque propio de una startup}.

\begin{importantbox}{«El 90\% de los avances vienen de nosotros»}
\textit{«El 90\% de los avances que sustentan la IA moderna fueron logrados por Google Brain, Research o DeepMind.» («90\% of breakthroughs that underpin modern AI were done by Google Brain, Research, or DeepMind.»)}

Esto no es solo autoelogio. La arquitectura transformer (Google Brain), BERT, T5 (Google Research), AlphaGo, AlphaFold (DeepMind): todo esto constituye, en efecto, la infraestructura de la IA moderna. El objetivo de la integración era concentrar estas fuerzas de innovación dispersas, evitando la competencia interna redundante y el desperdicio de recursos.
\end{importantbox}

\subsection{La estrategia de modelo de tres capas}

Tras la integración, Google DeepMind adoptó una estrategia clara de modelo de tres capas:

\begin{table}[H]
\centering
\begin{tabular}{p{0.15\textwidth} p{0.2\textwidth} p{0.5\textwidth}}
\toprule
\textbf{Capa} & \textbf{Producto representativo} & \textbf{Posicionamiento y estrategia} \\
\midrule
Modelo de frontera & Serie Gemini & Máximo rendimiento, código cerrado, producto comercial central \\
Modelo de código abierto & Serie Gemma & Se publica con un rezago de aproximadamente 6 meses respecto de la frontera; sirve al ecosistema de desarrolladores y a la comunidad académica \\
Modelo especializado por dominio & AlphaFold, entre otros & Optimizado a fondo para dominios científicos específicos, crea un valor insustituible \\
\bottomrule
\end{tabular}
\caption{La estrategia de modelo de tres capas de Google DeepMind}
\end{table}

\begin{knowledgebox}{Avances técnicos recientes de DeepMind}
Además del modelo central Gemini, DeepMind ha logrado avances importantes en varias direcciones recientemente:
\begin{itemize}
    \item \textbf{Modelos de video}: comprensión y creación de video generativo
    \item \textbf{Modelos interactivos del mundo}: capaces de simular y predecir los cambios del mundo físico
    \item \textbf{Serie de código abierto Gemma}: ofrece a la comunidad de desarrolladores modelos base de código abierto de alta calidad
\end{itemize}
Estas direcciones constituyen en conjunto la matriz técnica hacia la AGI: comprensión del lenguaje (Gemini), percepción visual (modelos de video), razonamiento físico (modelos del mundo) y construcción de ecosistema (Gemma).
\end{knowledgebox}

\subsection{La decisión estratégica de quedarse en el Reino Unido}

Un detalle interesante es que, aunque Silicon Valley es el centro global de la IA, DeepMind eligió \textbf{quedarse en el Reino Unido} para desarrollarse. Hassabis se muestra orgulloso de esto y sostiene que \textbf{es plenamente posible construir en Europa una empresa del orden del billón de dólares}.

La lógica detrás de esta decisión incluye: la profunda tradición académica del Reino Unido (la base de investigación en IA de Cambridge, Oxford y el UCL), un entorno regulatorio relativamente favorable, y la necesidad de mantener la independencia de la cultura organizacional.

\begin{knowledgebox}{Hassabis y Elon Musk}
Hassabis menciona su encuentro temprano con Elon Musk: los dos compartían una ambición común respecto de la AGI. Fue precisamente en su intercambio con Hassabis que Musk comprendió más a fondo el potencial y los riesgos de la AGI, lo cual impulsó indirectamente que Musk participara más tarde en la fundación de OpenAI y, después, de xAI. Aunque los dos tomaron caminos comerciales distintos, coinciden en gran medida en el juicio central de que «la AGI transformará por completo el mundo».
\end{knowledgebox}

\subsection{Resumen de la sección}

A través de la integración tres en uno, Google DeepMind centralizó fuerzas de innovación dispersas y adoptó una estrategia de modelo de tres capas: frontera, código abierto y especializado por dominio. La decisión de la organización de quedarse en el Reino Unido muestra la posibilidad de construir, fuera de Silicon Valley, un laboratorio de IA de nivel mundial.

%% ==================== Section 5 ====================
